\section{图像处理与深度学习}

\subsection{图像的矩阵表示}

\subsubsection{灰度图像}

对于一个图像，我们尝试用 $n \times m$ 的矩阵去表示它。对于矩阵每个位置 $(i, j)$ 上的元素，其值就是图像在这个位置的灰度值。这个位置的值被称为该位置的\textbf{采样}。

矩阵的大小越大，图像就越清晰，所谓\textbf{采样率}就是反映了采样有多精细。矩阵的元素值范围越大，图像的灰度就越丰富，所谓\textbf{量化等级}就是元素值的范围大小。

通过采样，图像在空间上和幅值上都从连续的量变成了离散的量。这个过程被称为图像的数字化，得到的图像也称为\textbf{离散图像}。

在数字图像处理中，一般都取图像的阵列大小 $N$ 和量化等级 $C$ 为 $2$ 的次幂。这是为了方便快速傅里叶变换的计算以及计算机存储的方便。

\subsubsection{彩色图像}

在彩色图像中，每个像素点由三个值表示，分别是红、绿、蓝三个分量。因此，彩色图像可以用三个灰度图像来表示。这种表示方式称为\textbf{RGB 表示}。

然而，人眼对于亮度的敏感度远大于对颜色的敏感度。因此还有\textbf{YUV 表示}：$Y$ 是两度，$U, V$ 是图像的色度在红、蓝方向的投影。使用 YUV 储存有利于图像压缩，因为很多图像的 $Y$ 分量饱和而 $U, V$ 分量稀疏，可以对于不同的分量矩阵采用不同的采样率（如 $4:2:2,\ 4:2:0$ 等）。处于 YUV 表示的特性，压缩 $U, V$ 分量对人眼影响不大，这个方法称为\textbf{降采样}。

\subsection{视频的矩阵表示}

视频就是在图像表示上再加上一个时间域。

\subsection{图像变化的表达式}

数字图像都是实数矩阵，且大多为 $N \times N$ 方阵，因此可以使用矩阵乘法表示图像变换：
$$
\matr{A}' = \matr{P} \matr{A} \matr{Q}
$$

\subsubsection{二维离散余弦变换}

二维离散余弦变换（DCT）的表达式定义如下：
$$
\matr{A}'_{u, v} = \frac{2}{N} \sum_{x = 0}^{N - 1} \sum_{y = 0}^{n - 1} \matr{A}_{x, y} \cos \frac{(2x + 1) u \pi}{2 N} \cos \frac{(2y + 1) v \pi}{2 N}
$$
